\chapter{设计和使用通用页面模板}\label{ch:templates}
\section{模板库、本文布局与默认设置（0.1.35）}
简单模板以一张纸的宽高和四向边距为布局，正文自动换行分页。字号、行距和段距属于排版样式，从齿轮的“正文与段落”设置。自定义模板通过起始页组和循环页组组织源布局；循环的是布局，实际页码连续递增。

主窗口\ui{模板库}用于管理可复用模板，显示缩略图、页面版式、页组顺序和字段声明。导入保存独立库内副本。\ui{保存修改}更新当前库内模板并选中最新修订，旧修订可在\ui{显示历史修订}中查看。编辑待入库的最近文件时，保存修改会写回原 JSON。\ui{另存为新模板}创建独立副本。移出条目可恢复，已保存的文档使用其内嵌模板。

排版右侧\ui{文档资源}用于更换本文字体、编辑本文布局、更换模板、填写字段和定位实际适配问题。更换模板先试排确认，保留正文与冻结资源；编辑菜单统一撤销重做。新文档默认仍在左下角设置配置，齿轮里的模板库只提供跳转入口。

\section{工作台与字段填写}
自定义工作台左侧维护起始、循环与未使用布局，中间画布显示正文编号及连接，支持选框、拖动移动和右下角缩放，右侧显示所选对象属性。共用布局的引用与独立副本明确区分。简单模板只呈现纸张、边距和预览。草稿自动收集有效输入，没有各页的“应用草稿”按钮；无效输入保留并提示，取消整个工作台不修改外部文档。确认按钮明确为“应用到当前文档”或“保存为新模板”。

字段表单只显示模板声明的字段，确定后写回单份 Markdown，保留未改内容和换行。重复定义及语法错误需回主编辑器定位处理。布局结构检查与真实内容试排分开，试排使用独立资源解析器，不改变当前选样。旧模板继续保留原设计器与源码块绑定，迁移为显式副本操作。

\section{有序正文框、固定字段与 PDF}
每页可有多个正文框，按明确编号连接。页面顺序和本页框编号共同决定唯一正文流，画布显示编号与连接关系，不猜空间阅读顺序。字段框引用稳定标识，多处引用重复显示，不能触发续页；页码使用独立自动框。固定框放不下时报告溢出，不默认缩字或丢弃内容。

PDF 支持按指定页码顺序一次导入多页背景。默认采用首个所选页的物理尺寸，可确认或保留原纸张；修改纸张时背景分别按宽高拉伸铺满，文本框保持毫米值。PDF 和图片均只作为参考，不自动描摹。预印纸只写内容；空白纸只输出明确的框线和内容。

\section{单份 Markdown 命名字段}
整份文稿只有一份源码。字段标识与显示名、具体框名分开；多个框可引用同一标识。字段块外全部进入连续正文，字段文字不在正文重复。示例：
\begin{verbatim}
:::field name
张三
:::

:::field student_id
20260001
:::

# 实验目的
这里开始是连续正文。
\end{verbatim}

重复定义、必填缺失或为空、未声明字段、缺少结束标记和嵌套字段块都阻止正式导出，并可定位源码。缺少结束标记不吞掉后续正文。代码围栏中的标记按普通代码处理。可选字段缺失时留空；字段支持已支持的 Markdown，但整体不得跨出固定框。

\section{从印刷纸样识别正文框}
导入图片或 PDF 后，可直接在底图上调整自由、横线与方格正文框，增删区域并核对行列。方格文字逐格放置，换行后按框和循环页顺序继续。保存模板后再次使用时沿用已确认的区域。

\section{结束、补齐与适配}
末尾补齐循环页组默认关闭，正文在哪一页结束就在哪一页停止。开启后只补完已经进入的当前循环组，补齐页仍绘制固定信息。空正文或仅字段时输出起始组一次；没有起始组则输出首循环布局，勾选补齐时输出完整首循环。已填写但没有实际框承接的内容明确报错。

正文跨框传递未消费内容与排版状态，按新宽度断行，不截断 Markdown 重解析。表格沿用原子行组和重复表头，公式和图片遵守原有尺寸规则。完整循环无法取得进展时停止并定位，页数保护与导出检查保留。

模板 v3 与原生项目 v3 显式版本化。旧版不自动迁移；复制旧几何时记录源版本和哈希，保留旧文件，原绑定、锚点和倍率需人工核对，不根据名字改写正文。以下各节记录旧模板的兼容操作。

\section{兼容旧模板：内容与绑定}
模板规定“在哪里写”，Markdown 规定“写什么”，绑定规定“哪一块内容进入哪个填写区”。同一模板可以用于不同文档；同一文档也可换用模板并重新绑定。学校报告、记录单或自定义表单只是不同配置，不是应用固定的栏目。

\par\noindent\begin{tabularx}{\linewidth}{@{}P{32mm}Y@{}}
\toprule
对象 & 用法\\\midrule
正文区 & 未绑定到固定区的正文按顺序流入，可自动续页。\\
固定填写区 & 姓名、记录编号、结论等完整内容块的指定区域，不能内部跨页。\\
锚点 & 定义首字横坐标与基线，帮助文字对齐预印下划线。\\
规则线 & 空白纸上需要绘制的框线或折线；与正文分开安排执行顺序。\\
底图 & PDF/图片转换来的纸面参考，不自动描摹为书写路径。\\
内容字段 & 为填写区附上稳定编号、显示名称、类型和必填规则，便于复用及 Agent 操作。\\
\bottomrule
\end{tabularx}

\section{新建一个中性模板}
\entry{设置 → 元素 → 页面模板… → 新建 / 设计模板…}
\begin{enumerate}
\item 确定真实纸张宽高。默认起点为普通 A4，不自动套用某所学校的表单。
\item 切换首页和续页，设置各自正文区。坐标原点在纸张左上角，横向向右、纵向向下，单位毫米。
\item 用\ui{绘制填写区}建立固定区；用\ui{放置锚点}设置首字基线；用\ui{绘制直线}添加需要书写的规则线。
\item 在属性表核对名称、位置和尺寸，修改数值后点击\ui{写入}。未写入的修改需先处理，再切换对象或确认。
\item 检查几何预览与问题列表，通过后导出模板 JSON，回到页面模板窗口绑定内容。
\end{enumerate}

当前每种页型只有一个正文流区域；首页和续页正文等宽，续页正文高度不少于首页。正文与固定区不能重叠。暂不支持多栏、不同栏宽的自动重排或任意页型序列。

\section{画布操作与精确调整}
\ui{选择 / 移动}工具可拖动区域、锚点或整条折线，区域四角可改尺寸，折线顶点可单独移动。每次拖动松开形成一次撤销操作，Esc 可取消本次拖动。多选后可使用对齐与等距分布；排列前核对选中的对象。

默认吸附间距为 1 mm，可修改或关闭，按住 Alt 可临时关闭吸附。标尺、缩放与底图透明度帮助观察，但不改变毫米坐标。拖动导致越界、重叠或违反页型约束时不会接受该结果；数值/JSON 草稿的错误通过问题列表修正。

给区域命名后再绑定内容。已在使用的名称需要修改时，优先使用外层页面模板窗口的\ui{联动重命名填写区 / 锚点…}，它会一起更新当前项目中的引用。直接改 JSON 名称不会自动推断原来的绑定关系。

\section{使用 PDF 或图片作参考}
设计器中的\ui{导入当前页底图…}支持本地 PNG/JPEG/PDF。PDF 导入时先选择页码，再确认目标是首页还是续页；导入 PDF 第 2 页与选择“续页”是两件事。

\begin{description}
\item[仅替换底图] 保留当前纸张和区域。若 PDF 尺寸不同，参考图会按当前纸张铺满，可能改变比例；需认真核对。
\item[按 PDF 本页尺寸新建布局] 按页面物理尺寸建立通用正文区，替换旧几何、底图和偏移。原有绑定需重新检查。
\end{description}
PDF 不会自动识别姓名栏、框线或正文。应使用底图手动设计，或让 Agent 提供经过验证的模板。导入可取消；超时可改用更简单的 PDF 或预先准备的图片。

\section{绑定固定字段：完整示例}
设想一张通用记录单：顶部留“记录编号”和“日期”，下面是自动续排的正文。先在首页画好两个固定区，将正文区放在它们下方。然后在源码中让编号、日期和正文各自成为独立段落。
\par\noindent\begin{minipage}{\linewidth}
\begin{Source}
20261014

2026.10.14

12.50

13.25
\end{Source}
\end{minipage}\par\medskip
\begin{enumerate}
\item 回到\ui{页面模板…}，选择记录编号填写区和第一段源码块。
\item 按需选择首页/所有页、倍率和基线锚点，点击\ui{绑定所选源码块}。
\item 将第二段绑定到日期区；后面未绑定的两段继续进入正文流。
\item 确认页面模板设置，等待实际字形排版，检查固定区是否越界。
\end{enumerate}
同一源码块不能同时绑定到两个固定区。解除绑定后，该内容回到正文；不会删除源码。删除或整体替换一个绑定块后，绑定可能失效，应重新选择块或解除旧绑定。

若要每页重复，首页和续页都需有同名填写区及所用锚点。固定区内容不能靠自动截断或压扁字形塞入；倍率也不会同比缩小表格的独立毫米尺寸。

\section{为复用模板声明内容字段}
在设计器的\ui{内容字段…}中新建字段，选择填写区，填写稳定编号和用户可读名称；设置内容类型为文字、公式、表格或任意完整块，并决定是否必填、在哪些页面出现。

使用已声明字段时，页面策略、倍率与锚点由模板控制，需要回到设计器修改。每份文档仍显式选择自己的源码块，程序不会因为字段名叫“姓名”就自动寻找文字。

添加语义字段会使用模板 v2；旧版应用可能无法读取，保留原模板副本。模板声明不保存某份报告的人名或正文。删除字段声明也不会自动删除原绑定，清理时应分别检查。

\section{预印纸、空白纸与试写}
\par\noindent\begin{tabularx}{\linewidth}{@{}P{26mm}Y Y@{}}
\toprule
模式 & 页面预览 & 机器输出\\\midrule
预印纸 & 显示纸面底图，或灰色参考框线。 & 只写内容，不重复画印刷底图与参考框线。\\
空白纸 & 显示配置的规则线与内容。 & 先画显式规则线，再写各固定区与正文；不复制扫描图上的印刷标题。\\
\bottomrule
\end{tabularx}

先用少量定位内容试写，测量整体横纵偏移，再在页面模板中设置毫米校准偏移。预印模式移动书写内容、保持纸面参考不动；空白模式规则线与内容一起移动。当前仅支持平移，不校正旋转或纸张缩放。

\section{查看适配问题}
\entry{查看问题 → 适配报告… / 修正模板…}
等待最新预览完成后打开适配报告，检查必填字段、绑定状态、缺字和溢出。选中问题可定位具体填写区或源码；也可另存 JSON 供 Agent 处理。

“模板几何有效”只说明区域合法，不代表正文能放下；“内容可导出”也不代表已完成设备校准。改动正文、字号、模板或资源后，要重新查看最新报告。使用\ui{普通页面}会移除模板并清除固定区绑定，内容返回普通正文流。


\section{连续正文流中的分页与换框}
多文本框连续正文流中，\code{\framebreak} 跳到下一个编号正文框；同页下一框不增加页码。\code{\newpage}、\code{\pagebreak}、\code{\clearpage} 跳到下一物理页，跳过当前页剩余正文框。普通 Markdown 页面也支持分页。

命令可独占一行或位于段内；连续换框逐个推进，空白页开头的分页不额外造页，尾随控制不额外创建空白输出页。列表标记随正文进入新框。固定字段、表格单元格和图片说明不执行分页或换框；普通页面及旧版固定区不执行换框。\code{\pagebreak[4]} 强制分页；强度 0--3 不强制推进。控制命令写在公式外，原子公式本身不拆分。内容由 AI 根据 Markdown 规范生成后复制或导入，不增加插入控制符的按钮。旧语法项目打开时显式迁移换框命令，保留原文件并另存。
